\chapter{Renderer Registry, Capabilities, and Evidence}
\label{ch:feature-matrix}

This appendix begins with the canonical public registry at the pinned source snapshot, then
retains a narrower capability grid from an earlier executed audit. The two cardinalities are
deliberately different: CNA exposes 46 public identities, while the five EasyGL profiles share
one factory, yielding 42 concrete implementation factories. An identity absent from the older
grid is simply unaudited by that grid; absence is not an unsupported-feature claim.

\begin{small}
\begin{longtable}{>{\raggedright\arraybackslash}p{3.3cm}>{\raggedright\arraybackslash}p{8.1cm}@{\hspace{0.6cm}}r}
\toprule
\textbf{Family / chapter} & \textbf{Canonical public selectors} & \textbf{Count} \\
\midrule
\endfirsthead
\toprule
\textbf{Family / chapter} & \textbf{Canonical public selectors} & \textbf{Count} \\
\midrule
\endhead
OpenGL fixed function & \texttt{OPENGLES1}, \texttt{OPENGL1} & 2 \\
OpenGL programmable and portable & \texttt{OPENGLES2}, \texttt{OPENGLES3},
  \texttt{OPENGL33}, \texttt{WEBGL1}, \texttt{WEBGL2}, \texttt{OPENGL2},
  \texttt{OPENGL4}, \texttt{PORTABLEGL} & 8 \\
Native modern GPU APIs & \texttt{VULKAN}, \texttt{WEBGPU}, \texttt{SDL\_GPU},
  \texttt{METAL} & 4 \\
Abstraction layers I & \texttt{BGFX}, \texttt{MAGNUM}, \texttt{LLGL},\newline
  \texttt{DILIGENT} & 4 \\
Abstraction layers II & \texttt{SOKOL}, \texttt{WICKED}, \texttt{FNA3D} & 3 \\
DirectX 1--3 and free-direct & \texttt{DIRECTX1}, \texttt{DIRECTX2},
  \texttt{DIRECTX3}, \texttt{FREEDIRECT} & 4 \\
Retro ladder & \texttt{DIRECTX5}, \texttt{DIRECTX6}, \texttt{DIRECTX7},
  \texttt{DIRECTX8}, \texttt{GLIDE} & 5 \\
Modern Direct3D & \texttt{DIRECTX9}, \texttt{DIRECTX10}, \texttt{DIRECTX11},
  \texttt{DIRECTX12} & 4 \\
Windows 2D & \texttt{DIRECT2D}, \texttt{GDI} & 2 \\
2D and vector rasterizers & \texttt{SDL\_RENDERER}, \texttt{SKIA},
  \texttt{BLEND2D}, \texttt{OPENVG} & 4 \\
Web and DOM & \texttt{CANVAS}, \texttt{HTML\_DOM},\newline \texttt{SVG\_DOM} & 3 \\
Non-GPU and diagnostic & \texttt{HEADLESS}, \texttt{SOFTWARE}, \texttt{STUB} & 3 \\
\midrule
\textbf{Total} & \textbf{46 public identities / 42 implementation factories} & \textbf{46} \\
\bottomrule
\end{longtable}
\end{small}

\texttt{ASCII} is not a current identity; its reusable logic moved to
\cnaclass{CNA::Graphics::AsciiPostProcessEffect}. \texttt{EASYGL}, \texttt{DX3},
\texttt{D3D9}, \texttt{D3D11}, and \texttt{D3D12} are not accepted selectors. The defaults
are \texttt{OPENGLES3} on Linux, \texttt{WEBGL2} under Emscripten, and
\texttt{SDL\_RENDERER} elsewhere. Chapter~\ref{ch:renderer-selection-identity} explains the
identity/factory distinction; \texttt{docs/renderer-registry.md} is the primary source.

\section{A capability grid, condensed from the project's own master matrix}

The registry above tells the current identity story; this section preserves the
complementary, narrower view from the earlier seven-renderer executed audit recorded in
\texttt{docs/graphics-renderer-\allowbreak{}feature-matrix.md}:
(\texttt{SDL\_RENDERER}, the EasyGL implementation, \texttt{VULKAN}, \texttt{BGFX}, and
the Direct3D~9/11/12 implementations). The historical D9/D11/D12 column labels below name
native APIs, not accepted current selectors. \texttt{WEBGPU} and \texttt{SDL\_GPU} were
excluded from that source document's own comparison, the former by its own documented choice
pending a broader feature surface, the latter because it postdates the document; see
Chapter~\ref{ch:webgpu-backend} and Chapter~\ref{ch:sdlgpu-backend} for their own coverage
instead. Status symbols match the source document's own legend, spelled out as plain text
rather than a symbol font this project's preamble doesn't load: \textbf{OK} correct and
verified, \texttt{X} known gap not fixed, \texttt{BLK} formally BLOCKED (needs a
project-owner architecture decision), \texttt{P} the compact table form of the source's
\texttt{PARTIAL} (implemented but not independently pixel-verified or only partially covered),
\texttt{N/A} not applicable to this renderer's own scope, \texttt{---} not attempted. Every cell
below is a condensed transcription of a real row
in the source document --- read that document directly (or the renderer's own chapter, cited in
the per-renderer table above) for the full evidence and task numbers behind each symbol; this
grid answers ``does it work'' at a glance, not ``how do we know.''

\begin{center}
\footnotesize
\setlength{\tabcolsep}{3pt}
\begin{longtable}{p{3.5cm}*{7}{p{1.35cm}}}
\toprule
\textbf{Feature} & \textbf{EGL} & \textbf{VK} & \textbf{BGFX} & \textbf{SDLR} & \textbf{D9} & \textbf{D11} & \textbf{D12} \\
\midrule
\endhead
\multicolumn{8}{l}{\textit{2D SpriteBatch / SpriteFont}} \\
All \texttt{Draw} overloads, sort modes, rotation/flip/scale/crop & OK & --- & --- &
  OK & OK & OK & OK \\
Custom \texttt{Effect} via \texttt{Begin(effect)} & OK & OK &
  X\textsuperscript{a} & N/A\textsuperscript{b} & OK & OK & OK \\
\texttt{SpriteFont} glyph placement/spacing/newline/fallback/flip & OK & --- & --- &
  OK & --- & OK & OK \\
\texttt{TextureAddressMode::}\allowbreak\texttt{Wrap} / \texttt{Mirror} via
  \texttt{SpriteBatch} & OK & OK & OK & BLK & OK & OK & OK \\
\midrule
\multicolumn{8}{l}{\textit{GraphicsDevice state objects}} \\
\texttt{BlendState} (all presets + custom factors/equations) & OK & OK &
  OK & OK & OK & OK & OK \\
\texttt{DepthStencilState} (compare func + full stencil ops) & OK & OK &
  P & OK & OK & OK & OK \\
\texttt{RasterizerState} aggregate\textsuperscript{d} & OK & P & P & OK & OK &
  OK & OK \\
Per-slot \texttt{SamplerState} (16 slots) & OK & OK & OK & OK &
  P & OK & OK \\
\texttt{GraphicsDevice.}\allowbreak\texttt{ReferenceStencil} & X & OK & X & N/A & --- & OK & N/A \\
\texttt{RenderTargetUsage} 2D color, Discard/P only\textsuperscript{e} &
  OK & OK & P & OK & P & P & P \\
\texttt{Clear} honors \texttt{ClearOptions::}\allowbreak\texttt{Stencil} &
  OK\textsuperscript{c} & P\textsuperscript{c} & P\textsuperscript{c} &
  N/A\textsuperscript{c} & P\textsuperscript{c} & OK\textsuperscript{c} &
  OK\textsuperscript{c} \\
\bottomrule
\end{longtable}
\end{center}

\begin{footnotesize}
\noindent\textsuperscript{a}\textbf{This cell deliberately corrects the source document.}
\path{docs/graphics-renderer-feature-matrix.md} still marks BGFX \textbf{OK}, but the live
renderer's \texttt{CompileProgram()} always fails and \cnaclass{ShaderEffect} exposes no
binary-loading route. Chapter~\ref{ch:shader-fx-gap} records the complete source audit.

\noindent\textsuperscript{b}\texttt{SDL\_RENDERER} rejects a non-null custom effect by design --- a
2D-only renderer with no programmable shader stage has no place for it to run. This is
\texttt{N/A}, not the source document's \texttt{BLK}: no unresolved project-owner architecture
decision is involved. See Chapter~\ref{ch:sdl-renderer}.

\noindent\textsuperscript{c}\textbf{This grid deliberately corrects the source document.}
\path{docs/graphics-renderer-feature-matrix.md} still marks this row \texttt{X} ``ignored, all
3 (Task~871, open)'' for \texttt{EasyGL} / \texttt{VULKAN} / \texttt{BGFX} --- stale as of this
book's own independent re-verification. Public stencil dispatch now exists. EasyGL has
selective D24S8 pixel proof. Vulkan/BGFX have cross-frame value proof but unconditionally clear
all view/pass attachments, hence \texttt{P}, not \textbf{OK}. SDL\_Renderer has no stencil
attachment and public masking makes the operation a no-op, hence \texttt{N/A}. D3D9's public
test proves colour preservation but not stencil bytes and gates target clears with the
backbuffer format, hence \texttt{P}; D3D11/D3D12 have direct GPU stencil-plane readback plus
the public dispatch route. Even their \textbf{OK} assumes \texttt{Depth24Stencil8}: CNA's shared
format mask wrongly retains stencil for depth-only target formats. The complete qualification
is \S\ref{sec:backend-clear-matrix}.

\noindent\textsuperscript{d}\textbf{Do not read this aggregate cell field-by-field.}
The source document collapses culling, fill, bias, multisampling, and scissor enable into one
status. The live audit found that SDL\_Renderer ignores the whole rasterizer object while
applying a non-empty scissor rectangle unconditionally, and D3D12's otherwise-real rasterizer
PSOs still force a full-target native scissor. Section
\ref{sec:backend-viewport-scissor-matrix} supersedes this row for viewport/scissor decisions.

\noindent\textsuperscript{e}\textbf{This row is deliberately narrow.}
It compares only non-MSAA \cnaclass{RenderTarget2D} color across separate
\texttt{DiscardContents}/\texttt{PreserveContents} bind cycles. SDL Renderer, EasyGL, and
Vulkan have discriminating pixels; BGFX has smoke only; D3D9/11/12 are source-proven here.
It says nothing about Platform, depth/stencil, cube targets, presentation/backbuffer usage,
redundant binds, or plural cube bindings. Those cases diverge materially; see
\S\ref{sec:backend-rendertarget-usage-matrix}.
\end{footnotesize}

This grid omits the source document's own Texture2D/Texture3D/TextureCube,
Model/OcclusionQuery, and remaining-genuine-limitations tables entirely --- not because they are
less real, but because condensing all six original tables into one appendix would defeat the
``quick reference, not replacement'' purpose this appendix states for itself; the two grids
above were chosen as the ones a reader deciding which renderer to target for a specific game
mechanic (2D-only? relies on device state objects for a real gameplay effect like an outline
shader via a custom \texttt{DepthStencilState}?) is most likely to need first. The full six-table
original remains the primary source for everything else, subject to the explicitly documented
stale cells this appendix has independently re-audited.

\section{Reading this table safely}

Two cautions from this book apply to every row above. First, ``verified'' means something
different for different renderers --- Chapter~\ref{ch:direct3d-backends}'s own D3D9 column is
explicitly the strictest bar in the project (pixel-for-pixel against a real running XNA~4.0
runtime), while a BGFX or CANVAS row's own verification is scoped more narrowly, and each
renderer chapter says so directly rather than leaving the reader to assume equivalence.
Second, this table is a snapshot. This book's own source material carried dated corrections in
both directions --- the Direct3D~9 renderer's own documentation and the project's cross-renderer
feature matrix genuinely disagreed about the render-target-as-texture fix's timing at one
point, and the feature matrix's own account was the more current one. Treat this appendix the
same way: a starting map, not a final word, to be checked against the project's own current
\path{docs/graphics-renderer-feature-matrix.md} before relying on any single row.
